\documentclass[a4paper,12pt]{article}
\usepackage[utf8]{inputenc}
\usepackage[T1]{fontenc}
\usepackage[french]{babel}
\usepackage{geometry}
\geometry{left=2.5cm,right=2.5cm,top=2.5cm,bottom=2.5cm}

\usepackage{lmodern}
\usepackage{xcolor}
\definecolor{titleblue}{RGB}{0,51,140}
\definecolor{TITLEBLUE}{RGB}{0,51,140}

\usepackage{hyperref}

\pagestyle{empty}

\begin{document}

\begin{flushleft}
    \textbf{Ismail AISSAOUI} \\
    Ingénieur d'État en Intelligence Artificielle \\
    Chlef, Algérie \\
    ismail.aissaoui.pro@gmail.com \\
    +213 660 70 77 96 \\
    \href{https://ismail-aissaoui.vercel.app}{ismail-aissaoui.vercel.app}
\end{flushleft}

\begin{flushright}
    À l'attention du \\
    \textbf{Professeur Abdelaziz BOURAS} \\
    Laboratoire DISP \\
    Université Lumière Lyon 2 \\
    69621 Villeurbanne, France \\
    \medskip
    Le 29 Avril 2026
\end{flushright}

\bigskip

\noindent \textbf{Objet : Candidature pour un doctorat en Gestion du Cycle de Vie des Données (PLM) \& IA}

\bigskip

Monsieur le Professeur,

Reconnu pour vos travaux fondateurs sur le Product Lifecycle Management (PLM) et l'interopérabilité des systèmes, c'est un honneur pour moi de vous soumettre ma candidature pour un doctorat au laboratoire DISP. Mon profil d'ingénieur IA spécialisé dans les systèmes industriels s'inscrit directement dans vos thématiques de recherche sur la valorisation des données tout au long du cycle de vie.

Mon expérience actuelle à la Sonatrach porte sur la conception d'un \textbf{Jumeau Numérique (Architecture Générative Mamba-SSM)} pour turbines à gaz. J'y traite l'ensemble du pipeline de données, de l'acquisition haute fréquence sur le "Edge" (\textbf{C++ JIT xLSTM}) à la modélisation sémantique et prédictive. Ce projet m'a sensibilisé aux enjeux de continuité numérique et de gestion de la connaissance technique au sein d'infrastructures critiques.

Je souhaite explorer comment l'IA générative et les modèles de substituts peuvent transformer la gestion du cycle de vie des produits, en permettant des simulations de dégradation beaucoup plus précises et actionnables. Votre vision du cycle de vie des systèmes est le socle sur lequel je souhaite bâtir mes travaux de recherche.

Je vous joins mon CV et reste à votre disposition pour un échange approfondi sur mes propositions de recherche.

Je vous prie d'agréer, Monsieur le Professeur, l'expression de mes salutations les plus distinguées.

\bigskip

\begin{flushright}
    \textbf{Ismail AISSAOUI}
\end{flushright}

\end{document}
